%&latex
\documentclass[a4paper]{article}

\frenchspacing

\usepackage[cp1250]{inputenc}
\usepackage[czech]{babel}

\usepackage{a4wide}
\usepackage{amsmath, amsthm, amssymb, amsfonts}
\usepackage[mathcal]{eucal}
\usepackage{graphicx}
\usepackage{url}
\usepackage{color}
\usepackage{wrapfig}
\usepackage{capt-of}
\usepackage{float}
\usepackage{ifthen}
\usepackage{subfig}
\usepackage[normalem]{ulem}


% sirka a vyska textu nastavena jako papir, vsechny okraje vynulovany a pridano 20pt na kazdou stranu
% horizontalni rozmery
\setlength{\textwidth}{\paperwidth}
\addtolength{\textwidth}{-40pt}
\addtolength{\hoffset}{-1in}
\addtolength{\hoffset}{20pt}
\setlength{\oddsidemargin}{0in}
\setlength{\marginparsep}{0in}
% vertikalni rozmery
\setlength{\textheight}{\paperheight}
\addtolength{\textheight}{-60pt}
\addtolength{\voffset}{-1in}
\addtolength{\voffset}{20pt}
\setlength{\topmargin}{0in}
\setlength{\headheight}{0in}
\setlength{\headsep}{0in}


%Obrazek na miste
%pouziti
%%\obrazeknahore{adresa}{popisek}{label}
\long\def\obrazeknahore#1#2#3 {

\begin{figure}[t]
    \centering
    \includegraphics[width=0.8\textwidth]{#1}
    
    \caption{#2}
    \label{#3}
    
\end{figure}

}


%==========================================
%PEKELNA MAKRA NA ZAROVNANI OBRAZKU DOPRAVA

\makeatletter


%tohle je makro, ktere mi dovoluje obtekani i u kratkych environmentu
%ABSOLUTNE nechapu, jak to funguje, ale funguje to
%viz http://tex.stackexchange.com/questions/26078/ 
\def\odrovnej{\@@par
\ifnum\@@parshape=\z@ \let\WF@pspars\@empty \fi % reset `parshape'
\global\advance\c@WF@wrappedlines-\prevgraf \prevgraf\z@
\ifnum\c@WF@wrappedlines<\tw@ \WF@finale \fi}

\makeatother



%---
%makro, co da obrazek doprava a ostatni text ho obteka
%(bez toho predchazejiciho makra to ale poradne nebeha)
%pouziti:
%\obrazekvpravo{adresa}{popisek}{label}{procento sirky}
\long\def\obrazekvpravo#1#2#3#4{

\setlength\intextsep{-20pt}

    \begin{wrapfigure}{r}{#4\textwidth}
      \begin{center}
          \vspace{-10pt}
          
        \includegraphics[width=#4\textwidth]{#1}
        \vspace{-10pt}
        
      \end{center}
      
      \caption{#2}
      \label{#3}
      
      
    \end{wrapfigure}

\setlength\intextsep{0pt}

    
}




%---
%makro pro pripady, kdy wrapfigure neco mrsi
%je to docela pekelne
%je nutne mu dat jak text vpravo, tak text vlevo
%a nevim, jestli bude 100% fungovat, ale doufam, ze jo

%pouziti:
%\obrazekvpravominipage{adresa}{popisek}{label}{procento sirky}{1 - procento sirky}{text vlevo}
\long\def\obrazekvpravominipage#1#2#3#4#5#6{

\noindent\begin{minipage}{#5\linewidth}
\vspace{0pt}
#6
\end{minipage}
\hspace{0.5cm}
\noindent\begin{minipage}{#4\linewidth}
\vspace{0pt}
\centering
\includegraphics[width=0.9\textwidth]{#1}
\captionof{figure}{#2}
\label{#3}
\end{minipage}

}

%KONEC PEKELNYCH MAKER
%=====================

\def\lebka{\includegraphics[width=3mm]{../common/lebka}}

% makra pro poznamku u vyrokove a predikatove logiky
\def\vl{ -- ve v�rokov� logice}
\def\pl{ -- v predik�tov� logice}


%Vacsina prostredi je dvojjazicne. V pripade, ze znenie napr pozorovania je pisane po slovensky, malo by byt po slovensky aj oznacenie.

\def\ifis#1#2{\ifthenelse{\equal{#1}{0}}{}{#2}}

%\newenvironment{e}[3]{\pagebreak[2]\noindent\textbf{\ifis{#2}{$\bigstar$}#1}\ifis{#3}{\emph{~(#3)}}\par\noindent\leftskip 10pt }{\odrovnej\par\bigskip}
\newenvironment{e}[3]{\pagebreak[2]\noindent\textbf{\ifis{#2}{\lebka~}#1}\ifis{#3}{\emph{~(#3)}}\ifis{#2}{~\lebka}\par\noindent\leftskip 10pt }{\odrovnej\par\bigskip}


%obecne prostredie, ktore ma vyuzitie pri specialnych odstavcoch ako (uloha, algoritmus...) aby nevzniklo dalsich x prostredi
\newenvironment{obecne}[1]{\pagebreak[2]\noindent\textbf{#1}\par\noindent\leftskip 10pt}{\odrovnej\par\bigskip}

\newenvironment{report}{\pagebreak[2]\noindent\textbf{Report}\em\par\noindent\leftskip 10pt}{\par\bigskip}

%\newenvironment{reportN}[1]{\pagebreak[2]\noindent\textbf{Report~}\emph{(#1)}\emph\par\noindent\leftskip 10pt}{\odrovnej\par\bigskip}
\newenvironment{reportN}[1]{\pagebreak[2]\noindent\textbf{Report~}\emph{(#1)}\em\par\noindent\leftskip 10pt}{\odrovnej\par\bigskip}

\newenvironment{penumerate}{
\begin{enumerate}
  \setlength{\itemsep}{1pt}
  \setlength{\parskip}{0pt}
  \setlength{\parsep}{0pt}
  %\setlength{\topsep}{200pt}
  \setlength{\partopsep}{200pt}
}{\end{enumerate}}

\def\pismenka{\numberedlistdepth=2} %pouzit, ked clovek chce opismenkovany zoznam...

\newenvironment{pitemize}{
\begin{itemize}
  \setlength{\itemsep}{1pt}
  \setlength{\parskip}{0pt}
  \setlength{\parsep}{0pt}
}{\end{itemize}}

\definecolor{gris}{gray}{0.95}
\newcommand{\ramcek}[2]{\begin{center}\fcolorbox{white}{gris}{\parbox{#1}{#2}}\end{center}\par}


\title{\LARGE U�ebn� texty k st�tn� bakal��sk� zkou�ce \\ Matematika \\ Teorie graf�}

\begin{document}

\maketitle

\newpage
\setcounter{section}{16}
\setcounter{page}{109}
\section{Teorie graf�}

\begin{e}{Po�adavky}{0}{0}
\begin{pitemize}
    \item Z�kladn� pojmy teorie graf�, reprezentace grafu.
    \item Stromy a jejich z�kladn� vlastnosti, kostra grafu.
    \item Eulerovsk� a hamiltonovsk� grafy.
    \item Rovinn� grafy, barven� graf�.
\end{pitemize}
\end{e}

\subsection{Z�kladn� pojmy teorie graf�, reprezentace grafu}

\begin{e}{Definice}{0}{Graf}
Graf $G$ je uspo��dan� dvojice $(V,E)$, kde $V$ je n�jak� mno�ina a $E$ je mno�ina dvouprvkov�ch podmno�in mno�iny $V$ (tak�e neuspo��dan�ch dvojic). Prvky mno�iny $V$ se jmenuj� \textbf{vrcholy grafu $G$} a prvky mno�iny $E$ \textbf{hrany grafu $G.$}
\end{e}

\begin{e}{Definice}{0}{Orientovan� graf}
$G$ je dvojice $(V,E),$ kde $E$ je podmno�ina kart�zsk�ho sou�inu $V \times V$. Prvky $E$ (tj. uspo��dan� dvojice prvk� z $V$) naz�v�me orientovan� hrany grafu.
\end{e}

\begin{e}{Definice}{0}{Symetrizace grafu}
Orientovan�mu grafu $G=(V,E)$ p�i�ad�me neorientovan� graf $\mathrm{sym}(G) = (V,\bar E) $ kde $\bar E=\{\{x,y\};$ $(x,y) \in E \lor (y,x) \in E \}.$ Graf $\mathrm{sym}(G)$ se naz�v� \textbf{symetrizace} grafu $G.$ \textit{(Z orientovan�ho grafu se odstran� �daje o sm�ru hran.)}
\end{e}

\begin{e}{Definice}{0}{D�le�it� typy graf�}
\begin{pitemize}
\item \textbf{�pln� graf $K_n$:} $V=\{1,\dots,n\}, E=\binom{V}{2}$ \textit{(Ka�d� vrchol je spojen hranou s ka�d�m nap�. $K_5$ --- \uv{pentagram}.)}
\item \textbf{kru�nice $C_n$:} $V=\{1,\dots,n\}, E=\{\{i, i+1\}; i=1,\dots,n-1\} \cup\{\{1,n\}\}$ \textit{(V~kru�nici se nesm� opakovat vrcholy.)}
\item \textbf{cesta $P_n$:} $V=\{0,1,\dots,n\}, E=\{\{i-1,i\}; i=1,\dots,n\}$ \textit{(Jako kru�nice, ale bez posledn� hrany.)}
\item \textbf{bipartitn� graf:} $V=\{u_1,\dots,u_n\} \; \cup \; \{v_1,\dots,v_m\}$, $E \subseteq \{\{u_i,v_j\}; i=1,\dots,n, j=1,\dots,m\}$; v \textbf{�pln�m bipartitn�m grafu} je $E = \{\{u_i,v_j\}; i=1,\dots,n, j=1,\dots,m\}$ \textit{(Ka�d� vrchol z jedn� partity je spojen hranou pouze s n�kter�mi (v �pln�m z ka�d�m) vrcholem druh� partity. Nap�. $K_{3,3}$, �pln� bipartitn� graf na 3 a 3 vrcholech.)}
\end{pitemize}
\end{e}

\begin{e}{Definice}{0}{Sled, tah}
Pro graf $G=(V,E)$ definujeme \textbf{sled} jako posloupnost $(v_0,e_1,v_1,\dots,e_n,v_n),$ kde $e_i=\{v_{i-1},v_i\} \in E$ pro $i=1,\dots,n.$ \textbf{Tah} je sled, ve kter�m se ��dn� hrana neopakuje.
\end{e}

\begin{e}{Definice}{0}{Isomorfismus graf�}
Dva grafy $G, G'$ pova�ujeme za \textbf{isomorfn�} (zna��me $G\simeq G'$), pokud se li�� jen v ozna�en� vrchol� a hran, tj. pokud existuje vz�jemn� jednozna�n� zobrazen� $f: G \rightarrow G'$ tak, �e plat� $\{x,y\}\in E \Leftrightarrow \{f(x),f(y)\} \in E'.$
\end{e}

\begin{e}{Definice}{0}{Podgraf, indukovan� podgraf}
Graf $G$ je \textbf{podgrafem} grafu $G',$ jestli�e $V(G)\subseteq V(G')$ a $E(G)\subseteq E(G') \cap \binom{V(G)}{2}$. Pro \textbf{indukovan� podgraf} $G$ grafu $G'$ plat�, �e $V(G)\subseteq V(G')$ a $E(G) = E(G') \cap \binom{V(G)}{2}.$ \textit{(Indukovan� podgraf dostaneme vymaz�n�m n�kter�ch vrchol� p�vodn�ho grafu a v�ech hran tyto vrcholy obsahuj�c�ch.)}
\end{e}

\begin{e}{Definice}{0}{Souvislost}
Neorientovan� graf je \textbf{souvisl�}, jestli�e mezi ka�d�mi jeho dv�ma vrcholy v n�m existuje cesta. Pro orientovan� grafy definujeme \textbf{slabou souvislost} --- po symetrizaci se z n�j stane souvisl� neorientovan� graf --- a \textbf{silnou souvislost} --- ka�d� dva vrcholy lze spojit orientovanou cestou v obou sm�rech.
\end{e}

\begin{e}{Pozn�mka}{0}{0}
Pro orientovan� grafy nen� samotn� pojem \uv{souvislost} definov�n.
\end{e}

\begin{e}{Definice}{0}{$n$-souvislost}
Graf je vrcholov� \textbf{$n$-souvisl�}, pokud m� alespo� $n+1$ vrchol� a po odebr�n� libovoln�ch $n-1$ vrchol� dostaneme v�dy souvisl� graf. Podobn� (p�es odeb�r�n� hran) definujeme \textbf{hranovou $n$-souvislost.}
\end{e}

\begin{e}{Pozn�mka}{0}{0}
Vrcholov� $n$-souvislost je siln�j�� podm�nka ne� hranov� $n$-souvislost, proto�e p�i odeb�r�n� $n-1$ vrchol� m��eme (a v�t�inou mus�me) odebrat v�ce ne� $n-1$ hran.
\end{e}

\begin{e}{Definice}{0}{Komponenta souvislosti}
\textbf{Komponenta souvislosti grafu} je jeho maxim�ln� souvisl� podgraf. (Sjednocen�m v�ech komponent grafu dostaneme p�vodn� graf).
\end{e}

\begin{e}{Definice}{0}{Most}
\textbf{Most} je hrana, kter� nele�� ve 2-souvisl�m podgrafu (po jej�m odstran�n� se zv�t�� po�et komponent).
\end{e}

\begin{e}{Definice}{0}{Blok}
Blok je maxim�ln� vrcholov� 2-souvisl� podgraf grafu. Samotn� graf o~2 vrcholech spojen�ch jednou hranou je tak� blok. \textit{(2-souvisl� podgraf, ke kter�mu se ned� p�idat ��dn� vrchol, proto�e by p�estal b�t 2-souvisl�.)}
\end{e}

\begin{e}{Definice}{0}{Artikulace}
Artikulace je vrchol $v$ souvisl�ho grafu $G$ takov�, �e $G - v$ nen� souvisl�.
\end{e}
%\input blok.pic

\begin{e}{Definice}{0}{Blokov� graf}
je graf incidence (sousednosti) blok� a artikulac� --- artikulac�m a blok�m odpov�daj� vrcholy, hrany zna�� incidenci blok� a artikulac�.
\end{e}

\begin{e}{V�ta}{0}{0}
Blokov� graf souvisl�ho grafu je strom.
\end{e}

\begin{e}{Definice}{0}{Hranov� pokryt�}
Mno�inu $C \subseteq E$ v grafu $G=(V,E)$ nazveme \textbf{hranov�m pokryt�m,} pokud ka�d� vrchol $v \in V$ je obsa�en alespo� v jedn� hran� $e \in C.$
\end{e}

\begin{e}{Definice}{0}{Stupe� vrcholu}
\textbf{Stupe� vrcholu,} $\deg_{G}(v),$ je po�et hran grafu $G$ obsahuj�c�ch vrchol $v.$ V p��pad� orientovan�ho grafu rozli�ujeme \textbf{vstupn� stupe� vrcholu,} $\deg^+_G(v),$ co� je po�et vstupn�ch hran, podobn� \textbf{v�stupn� stupe� vrcholu.}
\end{e}

\begin{e}{Definice}{0}{P�rov�n�}
Ka�d� mno�ina vz�jemn� disjunktn�ch hran v grafu se naz�v� \textbf{p�rov�n�}.
\end{e}

\begin{e}{Definice}{0}{k-faktor grafu}
\textbf{k-faktor} je podgraf $G'=(V,E')$ grafu $G=(V,E)$ takov�, �e $(\forall v \in V) \deg_{G'}(v)=k.$
\end{e}

\begin{e}{Pozn�mka}{0}{1-faktor grafu}
\textbf{1-faktor} je vlastn� p�rov�n� na v�ech vrcholech (�pln� p�rov�n�).
\end{e}

\begin{e}{V�ta}{0}{Princip sudosti}
$\sum_{v \in V} \deg_G(v) = 2 |E(G)|$ \end{e}

\begin{e}{Definice}{0}{Sk�re grafu}
\textbf{Sk�re grafu} je posloupnost stup�� vrchol� grafu, p�i�em� nez�le�� na po�ad�, v jak�m jsou uv�d�ny.
\end{e}

\begin{e}{V�ta}{0}{V�ta o sk�re}
Nech� $D=(d_1,\dots,d_n)$ je posloupnost p�irozen�ch ��sel. P�edpokl�dejme, �e $d_1\leq d_2 \leq \dots \leq d_n$ a ozna�me symbolem $D'$ posloupnost $(d'_1, \dots, d'_{n-1}),$ kde
$$
d'_i =
\begin{cases}
d_i & \text{pro $i < n-d_n$}\\
d_i-1 & \text{pro $i \geq n-d_n$}
\end{cases}
$$
Potom $D$ je sk�re grafu, pr�v� kdy� je $D'$ sk�re grafu. \textit{(Jakoby odebereme posledn� vrchol ($V_n$) a mysl�me si, �e byl propojen s p�edchoz�mi $d_n$ vrcholy.)}

\medskip
\begin{e}{D�kaz}{0}{0}
Jedna implikace je trivi�ln�, druh� (m�me $D$ sk�re grafu -- a tedy k n�mu n�jak� graf $G$ a dokazujeme, �e $D'$ je taky sk�re grafu $G'$) nen� o moc t잚� -- \uv{p�epoj�me} hrany tak, aby z vrcholu $v_n$ �ly pr�v� do $v_{n-d_n},\dots v_{n-1}$ a $v_n$ odebereme.
\end{e}
\end{e}


\begin{e}{Definice}{0}{Metrika grafu}
Pro souvisl� graf $G$ definujeme \textbf{metriku} jako funkci $d_G: V \times V \rightarrow \textbf{R},$ kde ��slo $d_G(v,v')$ p�edstavuje d�lku nejkrat�� cesty z $v$ do $v'.$ Funkce $d_G$ m� n�sleduj�c� vlastnosti (tj. spl�uje axiomy metriky, jak ji zn�me z metrick�ch prostor�):
\begin{penumerate}
\item $d_G(v,v') \geq 0,$ a $d_G(v,v')=0 \Leftrightarrow v=v';$
\item $(\forall v, v' \in V)(d_G(v,v')=d_G(v',v))$ \hfill \textit{(symetrie)}
\item $(\forall v, v', v'' \in V)(d_G(v, v'')\leq d_G(v,v')+d_G(v',v''))$ \hfill \textit{(troj�heln�kov� nerovnost)}
\end{penumerate}
\end{e}


\begin{e}{Definice}{0}{N�kter� grafov� operace}
Nech� $G=(V,E)$ je graf. Definujeme
\begin{pitemize}
\item \textbf{odebr�n� hrany:} $G-e=(V,E \setminus \{e\}),$ kde $e \in E$ je hrana grafu G
\item \textbf{p�id�n� nov� hrany:} $G+\bar e=(V,E \cup \{\bar e\}),$ kde $\bar e$ je dvojice vrchol�, kter� nen� hranou v $G$
\item \textbf{odebr�n� vrcholu:} $G-v=(V \setminus \{v\}, \{e \in E; v \not \in e\}),$ kde $v \in V$ (odebereme vrchol $v$ a v�echny hrany do n�j zasahuj�c�)
\item \textbf{d�len� hrany:} $G\%e=(V \cup \{z\},((E \setminus \{\{x,y\}\}) \cup \{\{x,z\},\{z,y\}\}),$ kde $\{x,y\} \in E$ je hrana a $z \not \in V$ je nov� vrchol (na hranu $\{x,y\}$ \uv{p�ikresl�me} nov� vrchol~$z$).
\end{pitemize}
�ekneme, �e graf $G'$ je \textbf{d�len�} grafu G, pokud $G'$ je isomorfn� grafu vytvo�en�mu z grafu $G$ postupn�m opakov�n�m operace d�len� hrany.
\end{e}


\subsection{Reprezentace grafu}

\begin{e}{Definice}{0}{Nakreslen�}
Graf lze reprezentovat nap�. jeho nakreslen�m -- lze si pod t�m p�edstavit i grafick� zn�zorn�n� na pap�r. Form�ln� definici nakreslen� provedeme v sekci o rovinn�ch grafech.
\end{e}

\begin{e}{Definice}{0}{Matice sousednosti}
M�jme graf $G=(V,E)$ s $n$ vrcholy $v_1,\dots,v_n.$ \textbf{Matice sousednosti} grafu $G$ je �tvercov� matice $A_G=(a_{ij})^n_{i,j=1}$ ��du $n$ definovan� p�edpisem
$$
a_{ij} =
\begin{cases}
1 & \text{pro $\{v_i,v_j\} \in E$}\\
0 & \text{jinak}
\end{cases}
$$
(Po umocn�n� matice sousednosti $A^k$ p�edstavuje ��slo $a_{ij}^{(k)}$ po�et sled� d�lky $k$ z vrcholu $v_i$ do vrcholu $v_j$ v grafu $G.$)
\end{e}

\begin{e}{Definice}{0}{Matice vzd�lenost�}
Pro grafy s ohodnocen�mi hranami lze zkonstruovat \textbf{matici vzd�lenost�} --- je to matice sousednosti, do kter� se v p��pad�, �e hrana existuje, ukl�d� m�sto jedni�ky jej� ohodnocen�.
\end{e}

\begin{e}{Definice}{0}{Laplaceova matice}
$$
(L_G)_{uv} =
\begin{cases}
\deg u & u=v\\
-1 & \{u,v\} \in E\\
0 & u \not= v \land \{u,v\} \not \in E
\end{cases}
$$
\textit{(Na hlavn� diagon�le je stupe� vrcholu, kde vede hrana, tam je -1, jinde 0. % predpokladal bych ze kdo si precetl predchozi sekci, vi o co jde :-) ... Stupe� vrcholu viz odst. \ref{stupen_vrcholu} 
)}

\begin{e}{Pozn�mka}{0}{0}
Laplaceovu matici lze pou��t mj. k v�po�tu po�tu koster grafu, jak uvid�me v n�sleduj�c� sekci.
\end{e}
\end{e}

\begin{e}{Definice}{0}{Matice incidence}
��dky matice odpov�daj� vrchol�m, sloupce odpov�daj� hran�m. Prvek matice se rovn� $-1$, pokud v tomto vrcholu za��n� dan� hrana, $+1$ pokud tam tato hrana kon��, $0$ jinak. Neorientovan� grafy maj� u obou vrchol� hrany hodnotu $+1$.
\end{e}

\begin{e}{Definice}{0}{Seznam soused�}
Pomoc� dvou pol�; v jednom ��sla v�ech n�sledn�k�, v druh�m poli indexy ur�uj�c�, kde za��n� sekvence soused� dan�ho vrcholu.
\end{e}

\begin{e}{Definice}{0}{Seznam hran}
Pole s prvky o dvou slo�k�ch, zaznamen�vaj� se do n�j hrany ve form� obou jejich vrchol�; pro orientovan� grafy lze stanovit, �e prvn� slo�ka bude reprezentovat po��te�n� a druh� koncov� vrchol hrany.
\end{e}

\pagebreak[3]
\subsection{Stromy a jejich z�kladn� vlastnosti}

\begin{e}{Definice}{0}{Strom, list}
\textbf{Strom} je souvisl� graf neobsahuj�c� kru�nici. \textbf{List} (koncov� vrchol) je vrchol stupn� 1.
\end{e}

\begin{e}{V�ta}{0}{0}
Po�et strom� na $n$-vrcholov� mno�in� je $n^{n-2}.$
\end{e}

\begin{e}{V�ta}{0}{Charakterizace strom�}
Pro graf $G=(V,E)$ jsou n�sleduj�c� podm�nky ekvivalentn�:
\begin{penumerate}
\item $G$ je strom.
\item Pro ka�d� dva vrcholy $x,y \in V$ existuje pr�v� jedna cesta z $x$ do $y.$ \textit{(jednozna�nost cesty)}
\item Graf $G$ je souvisl� a vynech�n�m libovoln� hrany vznikne nesouvisl� graf. \textit{(minim�ln� souvislost)}
\item Graf $G$ neobsahuje kru�nici a ka�d� graf vznikl� z $G$ p�id�n�m hrany ji� kru�nici obsahuje. \textit{(maxim�ln� graf bez kru�nic)}
\item $G$ je souvisl� a $|V|=|E|+1.$
\item $G$ je acyklick� a $|V|=|E|+1.$
\end{penumerate}
\end{e}

\begin{e}{Definice}{0}{Kostra grafu}
\textbf{Kostra grafu} $G$ je strom, kter� je podgrafem $G$ a obsahuje v�echny vrcholy grafu~$G$.
\end{e}

\begin{e}{V�ta}{0}{Po�et koster}
Po�et koster grafu $G$ je $\det(L'_G),$ kde $L'_{G}$
je matice, kter� vznikne odstran�n�m $i$-t�ho ��dku a $i$-t�ho sloupce z~Laplaceovy matice charakterizuj�c� graf.

??? Dukaz
\end{e}

\begin{e}{Report}{0}{0} Representace grafu byla docela jednoducha, popsal jsem ty representace, hlavni rozdil byl v casove narocnosti nalezeni sousedu vrcholu, bavili jsem se o BFS, DFS, jak jednotlive representace tedy ovlivni casovou narocnost. Popsal jsem a dokazal, co jsou mocniny matice sousednosti.
\end{e}
\begin{e}{Report}{0}{Majerech} teorie grafu: zakl. pojmy, souvislost, stromy (jen definice), kostry (chtel obecnou definici i pro nesouvisle grafy); 
eulerovy \& hamilt. grafy, rovinne grafy a barveni - definice a informace o tom, jak jsou problemy tezke z algoritmickeho hlediska. 
\end{e}
\begin{e}{Report}{0}{Fiala} Za�al jsem stromy a kostrami, prakticky to, co je v materi�lech + napsal jsem d�kaz Cayleyho formule (je opravdu lehk� a zabere hodn� pap�ru). M�l jsem je�t� dok�zat n�jakou implikaci u t�ch strom�. Tam jsem se tro�i�ku zamotal, ale nakonec +- 
\end{e}
\begin{e}{Report}{0}{Fiala} Stromy, min. kostra - Uvedl jsem definice, ekvivalentn� podm�nky, �e graf je stromem (bez d�kazu), Kruskal�v a Prim�v algoritmus. U Kruskalova alg. jsem se zamotal p�i vysv�tlov�n�, pro� je korektn�, o Primov� alg. ani slo�itosti obou m� nenechal mluvit (ale m�l jsem n�co na pap��e a �etl si to). 
\end{e}
\begin{e}{Report}{0}{Surynek}napsal jsem ekvivalentni definice stromu, par ekvivalenci dokazal, par obecnych kydu, definoval jsem kostru a vyslovil par vet pro kostry (pocet koster uplnaku, pocet koster grafu pres Laplaceovu matici), rekl jsem, jak se algoritmicky hleda kostra. Krom tech ekvivalenci u stromu jsem neumel nic dokazat (pocet koster uplnaku jsem tusil obrysy Obratlovc�, ale jinak nic moc, pr� �e je to "takov� n�jak� podivn�"), proto za 2. 
\end{e}
\begin{e}{Report}{0}{Ml�ek} Definice, z�kladn� tvrzen�. Nebylo toho moc, tak jsem pro jistotu p�ipsal je�t� n�co o minim�ln� kost�e a nast�nil algoritmy (Jarn�k, Bor�vka, Kruskal). U tvrzen� o po�tu koster �pln�ho grafu o (resp. po�tu strom� na) n vrcholech se m� ptal na d�kaz, nev�d�l jsem ho, ale to z�ejm� nevadilo (ostatn� bylo snad bez v�tek). P�r obecn�ch ot�zek kolem. 
\end{e}
\begin{e}{Report}{0}{Ml�ek} Na stromy jsem �ekl ideu d�kazu Cayleyho formule a fakt sem se divil, kdy� se m� po drsnostech z diferenci�ln�ch rovnic zeptal "No a v�te pros�mv�s, co je to indukovan� podgraf??" .. tak�e z trivialit prv�ku sem dostal jedni�ku, celkov� jedna. 
\end{e}
\begin{e}{Report}{0}{IOI 10.2.2011}
Teorie graf�: stromy
s)Uvedte alespo� dv� definice stromu a uka�te jejich ekvivalenci.
\\ 
b)Popi�te pojem kostra grafu a uve�te algoritmus pro jej� naleni v�etn� anal�zy pot�ebn�ho �asu a pam�ti
\end{e}
\subsection{Eulerovsk� a hamiltonovsk� grafy}

\begin{e}{Definice}{0}{Eulerovsk� graf}
Graf $G=(V,E)$ se naz�v� \textbf{eulerovsk�}, jestli�e existuje takov� po�ad� v�ech hran $e_1,\dots,e_n,$ �e $e_i \cap e_{i+1} \not = 0 \land e_1 \cap e_n \not = 0,$ tedy ka�d� dv� po sob� jdouc� hrany maj� spole�n� vrchol a rovn� prvn� a posledn� hrana se prot�naj�, ��dn� hrana se neopakuje. Jin�mi slovy: graf je eulerovsk�, pokud v n�m existuje uzav�en� sled $(v_0,e_1,v_1,\dots,e_{m-1},v_{m-1},e_m,v_0),$ v n�m� se ka�d� hrana vyskytuje pr�v� jednou. Takov� tah naz�v�me \textbf{uzav�en�m eulerovsk�m tahem.}
\end{e}

\begin{e}{V�ta}{0}{Charakteristika eulerovsk�ho grafu}
Graf $G=(V,E)$ je eulerovsk� pr�v� tehdy, kdy� je souvisl� a ka�d� vrchol $G$ m� sud� stupe�.

\medskip
\begin{e}{D�kaz}{0}{0}
\uv{$\Rightarrow$}: tato implikace je trivi�ln� --- eulerovsk� graf mus� b�t souvisl�, do ka�d�ho vrcholu mus�m vstoupit i z n�j vystoupit, co� zv��� stupe� v�dy o 2.

\bigskip
\noindent \uv{$\Leftarrow$}: M�jme tah maxim�ln� d�lky $v_0,e_1,\dots,e_m,v_m.$ Prvn� a posledn� vrchol tahu jsou toto�n�, jinak by do prvn�ho vrcholu vedl lich� po�et hran, a jeliko� graf m� v�echny stupn� sud�, dal by se tah prodou�it. Vid�me, �e ka�dou hranou v grafu vede n�jak� uzav�en� tah (nejdel�� tah hranou je uzav�en�, proto�e jinak by �el z vrcholu ze sud�m stupn�m prodlu�ovat). 

N� maxim�ln� tah obsahuje v�echny hrany a vrcholy, proto�e pokud ne, ze souvislosti jist� existuje vrchol, kter� je v max. tahu, z n�ho� vede hrana, kter� v max. tahu nen�. Tou vede uzav�en� tah a pokud ho spoj�me s na��m maxim�ln�m tahem, dostaneme del��, co� je spor.
\end{e}
\end{e}

\begin{e}{Definice}{0}{k-hamiltonovsk� graf}
Graf je \textbf{k-hamiltonovsk�}, pokud existuje posloupnost v�ech vrchol� $v_1,\dots, v_n$ takov�, �e $d_G(v_i,v_{i+1}) \leq k \land d_G(v_n,v_1) \leq k.$ \textit{(Do ka�d�ho vrcholu sm�me jen jednou.)} O grafu ��k�me jednodu�e, �e je \textbf{hamiltonovsk�}, pokud je 1-hamiltonovsk�.
\end{e}

\begin{e}{Definice}{0}{Hamiltonovsk� kru�nice}
Hamiltonovsk� kru�nice je kru�nice proch�zej�c� ka�d�m vrcholem grafu pr�v� jednou. Graf m� takovou kru�nici, pr�v� kdy� je hamiltonovsk�. Jej� nalezen� je NP-�pln� probl�m.

\begin{e}{Pozn�mka}{0}{0}
Probl�m nalezen� hamiltonovsk� kru�nice se d� up�esnit na probl�m nalezen� hamiltonovsk� kru�nice s nejmen�� vahou v ohodnocen�m grafu, co� je zn�m� probl�m obchodn�ho cestuj�c�ho. Je tedy tak� NP-�pln�, ale existuj� algoritmy, jejich� �e�en� jsou bl�zk� optim�ln�mu.
\end{e}
\end{e}


\begin{e}{V�ta}{0}{Diracova}
Graf $G=(V,E)$ je hamiltonovsk�, pokud plat�:
$$\forall v\in V: \deg v\geq \frac{|V|}{2}$$

\medskip
\begin{e}{D�kaz}{0}{0}
Dok�eme o n�co siln�j�� lemma pro jeden vrchol, z kter�ho Diracova v�ta plyne: pro dva vrcholy $u,v$ v grafu $G$ nespojen� hranou plat�, �e pokud $\deg u +\deg v \geq |V|$, potom je graf $G$ hamiltonovsk�, pr�v� kdy� $G$ s p�idanou hranou je hamiltonovsk�.

\medskip\noindent
Jedna implikace je trivi�ln�, tak�e vezmeme tu druhou. M�me v grafu $G+\{u,v\}$ hamiltonovskou kru�nici $C$. Pokud ta neobsahuje $\{u,v\}$ tak je i v grafu $G$ a kon��me, pokud tuto hranu obsahuje, ozna��me pro vrcholy grafu v po�ad�, v jak�m je proch�z� hamiltonovsk� kru�nice $u = v_1,v_2,\dots,v_n= v$:
\begin{align*}
A & :=\{i,\{u,v_i\}\in E\} \\
B & :=\{i + 1,\{v,v_i\}\in E\}
\end{align*}
Pokud maj� tyto mno�iny nepr�zdn� pr�nik, nalezneme vrcholy $v_k$ a $v_{k+1}$, p�es kter� m��eme kru�nici \uv{p�epojit}. $A$ ani $B$ neobsahuj� \uv{$1$}, tak�e $|A\cup B|\leq n-1$. Potom
$$|A\cap B|=|A|+|B|-|A\cup B|\geq |V|- |A\cup B|\geq 1$$
\noindent tak�e takov� bod $v_k$ existuje.
\end{e}
\end{e}

\begin{e}{V�ta}{0}{0}
Ka�d� souvisl� graf je 3-hamiltonovsk�.

\medskip
\begin{e}{D�kaz}{0}{0}
$G=(V,E)$ je souvisl�, existuje v n�m tedy kostra $T=(V,E'),$ kter� je souvisl�. Kostra vznikla ubr�n�m n�kter�ch hran, tak�e $d_T(x,y) \geq d_G(x,y) (\forall x,y \in V).$ Sta�� tedy dok�zat, �e kostra $T$ je 3-hamiltonovsk�.
\end{e}
\end{e}

\begin{e}{Lemma}{0}{0}
Ka�d� strom je 3-hamiltonovsk�.

\medskip
\begin{e}{D�kaz}{0}{0}
Indukc�:
\begin{penumerate}
\item pro $n \leq 4$ trivi�ln�
\item M�me dv� komponenty $T_1,T_2,$ ka�d� je 3-hamiltonovsk�. Graf je souvisl� $\rightarrow$ existuje most z $T_1$ do $T_2.$ Most vede p�es vrcholy $x',x,y,y',$ kde $x,x' \in T_1, y,y' \in T_2.$ Potom existuje 3-hamiltonovsk� propojen� komponent $T_1,T_2: x,(T_1),x',y',(T_2),y.$
\end{penumerate}
\end{e}
\end{e}

\begin{e}{V�ta}{0}{0}
Graf je vrcholov� 2-souvisl�, pr�v� kdy� v n�m mezi ka�d�mi dv�ma r�zn�mi vrcholy existuj� dv� vrcholov� disjunktn� cesty.

\medskip
\begin{e}{D�kaz}{0}{0}
Implikace $\Leftarrow$ je z�ejm�, opa�n� se d� uk�zat sporem -- nech� ve dvousouvisl�m grafu mezi n�jak�mi dv�ma vrcholy neexistuj� dv� vrcholov� disjunktn� cesty. Pak vezmu vrchol, kter� je na ka�d� cest� mezi t�mito dv�ma a odeberu ho -- a t�m se graf stane nesouvisl�m, co� je spor.
\end{e}
\end{e}

\begin{e}{V�ta}{0}{0}
Graf je vrcholov� 2-souvisl�, pr�v� kdy� jej lze vytvo�it z troj�heln�ku (t.j. z $K_3$) posloupnost� d�len� a p�id�v�n� hran (definice t�chto operac� jsou na za��tku kapitoly).
\end{e}

\begin{e}{V�ta}{0}{0}
Ka�d� 2-souvisl� graf je 2-hamiltonovsk�.

\medskip
\begin{e}{D�kaz}{0}{0}
??? Do ka�d�ho vrcholu vedou 2 vrcholov� i hranov� disjunktn� cesty --- p�i zp�te�n� cest� pou�iju druhou. D�kaz podobn� jako u v�ty o 3-hamiltonovsk�ch grafech.
\end{e}
\end{e}


\subsection{Rovinn� grafy}

\begin{e}{Definice}{0}{Oblouk}
\textbf{Oblouk} je podmno�ina roviny tvaru $o=\gamma(\langle0,1\rangle)=\{\gamma(x); x \in \langle0,1\rangle\}$, kde $\gamma:\langle0,1\rangle \to \mathbb{R}^2$ je n�jak� prost� spojit� zobrazen� intervalu $\langle0,1\rangle$ do roviny. P�itom body $\gamma(0)$ a $\gamma(1)$ naz�v�me \textbf{koncov� body} oblouku $o$.
\end{e}

\begin{e}{Definice}{0}{Nakreslen� grafu}
\textbf{Nakreslen�m grafu} $G=(V,E)$ rozum�me p�i�azen�, kter� ka�d�mu vrcholu v grafu $G$ p�i�azuje bod $b(v)$ roviny a ka�d� hran� $e=\{v,v'\}$ p�i�azuje oblouk $o(e)$ v rovin� s koncov�mi body $b(v)$ a $b(v').$ Zobrazen� $b(v)$ je prost� (r�zn�m vrchol�m odpov�daj� r�zn� body) a ��dn� z bod� tvaru $b(v)$ nen� nekoncov�m bodem ��dn�ho z oblouk� $o(e).$ Graf spolu s nakreslen�m naz�v�me \textbf{topologick� graf.}
\end{e}

\begin{e}{Definice}{0}{Rovinn� graf}
Nakreslen� grafu $G,$ v n�m� oblouky odpov�daj�c� r�zn�m hran�m maj� spole�n� nanejv�� koncov� body, se naz�v� \textbf{rovinn� nakreslen�}. Graf $G$ je \textbf{rovinn�,} m�-li alespo� jedno rovinn� nakreslen�.
\end{e}

\begin{e}{Definice}{0}{St�na grafu}
M�jme $G$ rovinn� topologick� graf. Mno�inu $A \subseteq \mathbb{R}^2\setminus X$ bod� roviny (kde $X$ je mno�ina v�ech bod� v�ech oblouk� nakreslen� grafu $G$) nazveme \textbf{souvislou,} pokud pro libovoln� dva body $x,y \in A$ existuje oblouk $o \subseteq A$ s koncov�mi body $x,y.$  Relace souvislosti rozd�l� mno�inu v�ech bod� roviny, kter� nele�� v ��dn�m z oblouk� nakreslen�, na t��dy ekvivalence. Ty naz�v�me \textbf{st�nami} topologick�ho rovinn�ho nakreslen� grafu $G$.
\end{e}

\begin{e}{Definice}{0}{Topologick� kru�nice}
Uzav�en� k�ivka v rovin� neprot�naj�c� sebe sama; form�ln� se definuje jako oblouk, jeho� koncov� body spl�vaj�.
\end{e}

\begin{e}{V�ta}{0}{Jordanova, o kru�nici}
Ka�d� topologick� kru�nice $k$ rozd�luje rovinu na pr�v� dv� souvisl� ��sti (\uv{vnit�ek} a \uv{vn�j�ek}), p�i�em� $k$ je jejich spole�nou hranic�.
\end{e}

\begin{e}{V�ta}{0}{Kuratowsk�ho}
Graf $G$ je rovinn�, pr�v� kdy� ��dn� jeho podgraf nen� isomorfn� d�len� grafu $K_{3,3}$ ani $K_5.$
\end{e}

\begin{e}{V�ta}{0}{Euler�v vzorec}
Nech� $G=(V,E)$ je souvisl� rovinn� graf, a nech� $s$ je po�et st�n n�jak�ho rovinn�ho nakreslen� $G.$ Potom plat� $|V|-|E|+s = 2.$

\medskip
\begin{e}{D�kaz}{0}{0}
Indukc� podle po�tu hran. Pro graf sest�vaj�c� pouze z jednoho vrcholu vzorec plat�.
\begin{penumerate}
\item Pokud p�idan� hrana nevytvo�� kru�nici, nezm�nil se po�et st�n, ale o jednu se zv�t�il po�et vrchol� a hran (graf mus� b�t v�dy souvisl�, tak�e jedin� mo�nost jak tohoto dos�hnout, je p�ipojit dal�� vrchol).
\item Pokud p�idan� hrana vytvo�� kru�nici, zv�t�� se po�et st�n o jednu (p�idan� hrana soused� se dv�ma r�zn�mi st�nami --- podle Jordanovy v�ty o kru�nici --- kter� p�ed p�id�n�m byly st�nou jedinou), po�et hran tak� o jednu, a po�et vrchol� se nezm�n�.
\end{penumerate}
Vzorec tedy v obou p��padech plat�.
\end{e}
\end{e}

\begin{e}{V�ta}{0}{2-souvisl� rovinn� graf}
Dvousouvisl� rovinn� graf m� hranice libovoln� st�ny libovoln�ho nakreslen� jako kru�nice.
\end{e}

\begin{e}{V�ta}{0}{Maxim�ln� po�et hran rovinn�ho grafu}
Nech� $G$ je rovinn� graf s alespo� 3 vrcholy. Potom $|E|\leq 3|V|-6.$ Rovnost nast�v� pro ka�d� maxim�ln� rovinn� graf, t.j. rovinn� graf, ke kter�mu nelze ji� p�idat ��dnou hranu (p�i zachov�n� mno�iny vrchol�) tak, aby z�stal rovinn�. Pokud graf $G$ nav�c neobsahuje troj�heln�k (t.j. $K_3$ jako podgraf) a m�-li alespo� 3 vrcholy, potom $|E|\leq 2|V|-4.$

\medskip
\begin{e}{D�kaz}{0}{0}
Ob� tvrzen� m��eme dok�zat pro maxim�ln� rovinn� graf. V prvn�m p��pad� je ur�it� dvousouvisl�, proto�e jinak m��u je�t� n�jak� dva body spojit. Nav�c ka�d� st�na mus� b�t troj�heln�k (je-li �tverec nebo n�co v�t��ho, tak� jdou je�t� n�jak� dva body spojit). T�m dostanu, �e $3s=2|E|$ a zbytek vyjde z Eulerova vzorce.

Druh� p��pad m� jednu zvl�tnost -- takov� graf m��e b�t hv�zda. Pro tu je ale tvrzen� spln�no. Pokud nen� hv�zda, u� mus� b�t dvousouvisl� a v�echny st�ny mus� b�t �tverce, tak�e dostanu $4s=2|E|$ a z Eulerova vzorce i cel� v�sledek.
\end{e}
\end{e}

\begin{e}{D�sledek}{0}{0}
Rovinn� graf m� alespo� jeden vrchol stupn� nejv��e 5.
\end{e}

\subsection{Barven� grafu}

\begin{e}{Definice}{0}{Neorientovan� graf s n�sobn�mi hranami (multigraf)}
je trojice $(V,E,\varepsilon),$ kde $V$ a $E$ jsou disjunktn� mno�iny a $\varepsilon: E \to \binom{V}{2} \cup V$ je zobrazen� \textit{(hrany prohl�s�me za \uv{abstraktn�} objekty a $\varepsilon$ ur�uje pro ka�dou hranu dvojici vrchol�, kter� jsou jej�mi \uv{konci}).}
\end{e}

\begin{e}{Pozn�mka}{0}{0}
Pro orientovan� graf je pojem definov�n obdobn�, pouze hrany jsou p�i�azeny uspo��dan� dvojici vrchol�, tedy: $\varepsilon: E \to \langle v_1,v_2 \rangle \cup V$.
\end{e}

\begin{e}{Definice}{0}{Geometrick� du�l grafu}
Nech� $G$ je topologick� rovinn� graf. Ozna�me $S$ mno�inu st�n $G.$ Jako \textbf{(geometrick�) du�l grafu} definujeme multigraf tvaru $(S,E,\varepsilon),$ kde $\varepsilon$ se definuje p�edpisem $\varepsilon(e)=\{S_i,S_j\},$ jestli�e hrana $e$ je spole�nou hranic� st�n $S_i$ a $S_j$ (p�i�em� m��e b�t $S_i=S_j,$ jestli�e z obou stran hrany $e$ je tat� st�na). Zna��me jej $G^*.$ \textit{(Sice opravdu plat� $G^{**}=G$, ale pak ji� pro spo��t�n� druh�ho du�lu pot�ebujeme zn�t du�l i~pro multigrafy!)}
\end{e}

\begin{obecne}{�loha \textit{(barven� mapy)}}
Uva�me politickou mapu, na n� jsou vyzna�eny hranice st�t�. P�edpokl�dejme, �e ka�d� st�t tvo�� souvislou oblast ohrani�enou n�jakou topologickou kru�nic�. Dv� oblasti pokl�d�me za sousedn�, jesli�e maj� spole�n� aspo� kousek hranice. Ka�d� st�t na takov� map� chceme vybarvit n�jakou barvou tak, �e sousedn� st�ty nikdy nebudou m�t stejnou barvu. Jak� minim�ln� po�et barev je pot�eba?
\end{obecne}

\begin{e}{Definice}{0}{Barevnost grafu}
Bu� $G=(V,E)$ graf, $k$ p�irozen� ��slo. Zobrazen� $b: V \rightarrow \{1,\dots,k\}$ nazveme \textbf{obarven�m grafu} $G$ pomoc� $k$ barev, pokud pro ka�dou hranu $\{x,y\}\in E$ plat� $b(x) \not = b(y).$ \textbf{Barevnost (chromatick� ��slo) grafu} $G,$ ozna�ovan� $\chi(G),$ je minim�ln� po�et barev pot�ebn� k obarven� $G.$
\end{e}

\begin{e}{Lemma}{0}{0}
Du�ln� graf rovinn�ho grafu je rovinn� graf.
\end{e}

\begin{obecne}{P�evod �lohy barven� mapy na �lohu hled�n� barevnosti grafu}
M�me-li mapu, kterou ch�peme jako nakreslen� n�jak�ho grafu $G,$ potom ot�zka obarvitelnosti mapy pomoc� $k$ barev je ekvivalentn� s obarvitelnost� du�ln�ho grafu $G^*$ pomoc� $k$ barev. Na druh� stran� plat�, �e ka�d� rovinn� graf se vyskytne jako podgraf du�ln�ho grafu n�jak�ho vhodn�ho grafu. Takto lze p�ev�st probl�m barven� map na probl�my barevnosti rovinn�ch graf�.
\end{obecne}

\begin{e}{V�ta}{0}{V�ta o p�ti barv�ch}
Pro ka�d� rovinn� graf $G$ plat� $\chi(G)\leq 5.$

\medskip
\begin{e}{D�kaz}{0}{0}
Indukc� dle po�tu vrchol� grafu. Pro $|V|\leq 5$ je tvrzen� trivi�ln�. Podle d�sledku v�ty o po�tu hran v rovinn�m grafu existuje vrchol stupn� $\leq 5.$ Pokud m� vrchol $v$ stupe� $< 5,$ pou�ijeme induk�n� p�edpoklad: obarv�me graf $G-v$ 5 barvami a $v$ p�i�ad�me barvu, kter� nen� mezi barvami jeho (nejv��e �ty�) soused�. Zb�v� tedy p��pad, kdy m� $v$ stupe� 5 a jeho sousedi jsou obarveni r�zn�mi barvami. Graf $G$ m� pevn� zvolen� rovinn� nakreslen� a v tomto nakreslen� budou $t,u,x,z,y$ soused� vrcholu $v$ v takov�m po�ad�, v jak�m p��slu�n� hrany vych�zej� z vrcholu $v$ (nap�. po sm�ru hodinov�ch ru�i�ek). Uva�me vrcholy $x$ a $y$ a nech� $V_{x,y}$ je mno�ina v�ech vrchol� grafu $G'=G-v$ obarven�ch barvami $b(x)$ nebo $b(y).$ Z�ejm� $x,y \in V_{x,y}.$

\medskip \noindent Nast�vaj� dva p��pady:
\begin{penumerate}
\item Neexistuje cesta z $x$ do $y$ pou��vaj�c� pouze vrchol� z $V_{x,y}.$
M�jme $V'_{x,y}$ mno�inu vrchol�, kter� jsou v $G'$ spojeny s $x$ cestou
pou��vaj�c� jen vrcholy z $V_{x,y}.$ Definujeme nov� obarven� $b'$ takto:
$b'(s) = b(s),$ pokud $s \not \in V'_{x,y};$ a pokud $s \in V'_{x,y},$
zm�n�me barvu $s$ z $b(y)$ na $b(x)$ nebo z $b(x)$ na $b(y)$ (tzn. na
mno�in� $V_{x,y}$ zam�n�me barvy). Z�ejm� $b'$ je obarven�, a proto�e
$b'(x)=b'(y)=b(y),$ m��eme polo�it $b'(v)=b(x).$ Tedy $b'$ je obarven�
grafu $G$ 5 barvami.
\item Pokud takov� cesta existuje (ozna�me ji $P$), uva�me vrcholy $t$ a $z.$
Z�ejm� $V_{x,y}$ a $V_{t,z}$ jsou disjunktn� mno�iny. Cesta $P$ spolu
s hranami $\{v,x\}$ a $\{v,y\}$ tvo�� kru�nici, kter� odd�luje $t$ a
$z,$ a proto by ka�d� cesta z $t$ do $z$ musela pou��t n�kter� vrchol t�to
kru�nice. Neexistuje tedy cesta z $t$ do $z$ pou��vaj�c� pouze vrchol� z
$V_{t,z}$ a obarven� grafu G 5 barvami lze zkonstruovat stejn� jako v
p�edchoz�m p��pad�, pouze mus�me za��t s vrcholy $z$ a $t.$
\end{penumerate}
\end{e}
\end{e}

\begin{e}{V�ta}{0}{Probl�m �ty� barev}
Je mo�n� ka�dou mapu obarvit 4 barvami?

\medskip
\begin{e}{D�kaz}{0}{0}
V�ta plat�, ale je velmi t�k� ji dok�zat (prob�r�n� mnoha p��pad� po��ta�em).
\end{e}
\end{e}

\begin{e}{Pozn�mka}{0}{0}
NP-�pln� probl�my: je d�n neorientovan� graf $G$ a ��slo $k.$
\begin{penumerate}
\item Je mo�n� $G$ obarvit $k$ barvami?
\item Tot�, ale p�edem v�me, �e k=3.
\item Tot�, ale graf je rovinn�.
\item Tot�, ale stupe� libovoln�ho vrcholu je nejv��e 4.
\item Je d�no obarven� t�emi barvami, dotaz na netrivi�ln� jin�.
\end{penumerate}
\end{e}


\subsection{Z�kladn� grafov� algoritmy}

\ramcek{10cm}{Tato sekce nen� po�adovan� ke zkou�ce!\\
\\
(Nebo teda je po�adovan�, ale v informatice, kde je vypracovan� zvl᚝)
}

V tomto odd�le zavedeme pro odhady �asov�ch slo�itost� algoritm� zna�en� $n=|V(G)|$ a $m=|E(G)|.$

\begin{e}{Algoritmus}{0}{Dijkstr�v algoritmus pro hled�n� nejkrat�� cesty}
M�me graf $G,$ jeho� hrany jsou ohodnoceny kladn�mi ��sly, tzn. �e je d�no zobrazen� $w: E(G) \to (0,\infty).$ (Ohodnocen� $w(e)$ hrany $e$ si p�edstavujeme jako jej� d�lku. D�lka cesty je rovna sou�tu d�lek jej�ch hran a vzd�lenost $d_{G,w}(u,v)$ vrchol� $u,v$ je rovna nejmen�� z d�lek v�ech cest spojuj�c�ch $u,v.$ \uv{Oby�ejn�} grafov� vzd�lenost je speci�ln� p��pad, toti� je-li $w(e)=1$ pro ka�dou hranu $e.$) Hled�me nejkrat�� cestu z vrcholu $s$ do v�ech ostatn�ch vrchol�.
\begin{penumerate}
\item (Inicializace) \hfil \\
    \uv{Odhad} vzd�lenosti $d(\cdot)$ u po��te�n�ho vrcholu $s$ nastav�me na 0,
    odhady u v�ech ostatn�ch vrchol� na $\infty$ (zn�me cestu d�lky 0 z $s$ do
    $s,$ d�lky ostatn�ch cest nezn�me). Do mno�iny $A$ vrchol�, u nich�
    je�t� nen� odhad definitivn�, d�me v�echny vrcholy krom� $s$.
\item (Volba mno�iny $N$) \hfil \\
    Do mno�iny $N$ pr�v� zpracov�van�ch vrchol� ulo��me v�echny vrcholy z $A,$ kter� maj� ze v�ech vrchol� z $A$ minim�ln� odhad vzd�lenosti; tyto vrcholy z $A$ vy�ad�me.
\item (Aktualizace odhad�) \hfil \\
    Pro ka�dou hranu $e=\{v,y\} \in E,$ kde $v \in N, y\in A,$ porovn�me hodnoty $d(y)$ a $d(v)+w({v,y}).$ Pokud $d(v)+w({v,y}) < d(y)$ (p�es vrchol $v$ vede do $y$ krat�� cesta, ne� jsme zat�m znali), nastav�me $d(y)$ na tuto hodnotu. Po vy�erp�n� v�ech takov�ch hran pokra�ujeme dal��m krokem.
\item (Test ukon�en�) \hfil \\
    Jestli�e odhady vzd�lenost� v�ech vrchol� v mno�in� $A$ jsou $\infty,$ algoritmus kon��, jinak pokra�uje krokem 2.
\end{penumerate}
Po skon�en� algoritmu je bu� $A = 0$ (je-li $G$ souvisl�) nebo $A$ obsahuje pouze vrcholy nedosa�iteln� cestou z vrcholu $s.$

??? Algoritmus jsem opravil, i kdy� je to vcelku zbyte�n� ... kdy�tak to n�kdo zkontrolujte. % -- Tuetschek
\end{e}

\begin{e}{V�ta}{0}{Slo�itost Dijkstrova algoritmu}
Lze ho implementovat v �ase $O(n \log n + m)$ --- nap�. pomoc� Fibonacciho hald.
\end{e}

\begin{e}{Pozn�mka}{0}{0}
Pokud hled�me nejkrat�� cestu pouze do jednoho zadan�ho vrcholu $c,$ m��eme ukon�it Dijkstr�v algoritmus hned pot�, co tento vrchol opust� mno�inu $A$ (jeho vzd�lenost se stane definitivn�). V�po�et tak� m��eme urychlit n�sleduj�c� heuristikou: m�me funkci $h: V(G) \to \langle0,\infty)$ spl�uj�c� $h(c)=0$ a pro ka�dou hranu $e=\{u,v\} \in E$ plat� $|h(u)-h(v)| \leq w(e)$ (v probl�mech dopravn�ho spojen� to m��e b�t nap�. vzd�lenost vzdu�nou �arou od c�le). Potom p�i volb� mno�iny $N$ vyb�r�me prvky s minim�ln�m sou�tem dosavadn�ho odhadu a heuristick� funkce --- $d(v)+h(v).$ Je-li $h$ kvalitn�, d� se �ekat, �e algoritmus najde definitivn� vzd�lenost do $c$ rychleji.
\end{e}


\begin{e}{Algoritmus}{0}{Prohled�v�n� do ���ky}
M�me graf $G=(V,E)$ a po��te�n� vrchol $s.$ Na za��tku polo��me $V_0 = \{s\},$ v dal��ch kroc�ch polo��me $V_{i+1} = \{v \in V \setminus (V_0 \cup \dots \cup V_i): \exists u \in V_i, \{u,v\} \in E\}.$ Slo�itost algoritmu je $O(n+m).$
\end{e}

\begin{e}{Algoritmus}{0}{Prohled�v�n� do hloubky}
??? Pom�rn� jasn�, �asto vede k exponenci�ln� slo�itosti. Nutno rozli�ovat metodu zpracov�n� prvk� --- \textbf{preorder, inorder} nebo \textbf{postorder.}
\end{e}

\begin{e}{Algoritmus}{0}{Hladov� (Kruskal�v) algoritmus na hled�n� minim�ln� kostry}
M�jme souvisl� graf $G=(V,E)$ s ohodnocen�m hran $w.$ Hrany m�me uspo��d�ny vzestupn� podle v�hy, $w(e_1) \leq \dots \leq w(e_n).$
\begin{penumerate}
\item Polo�me $E_0=0$
\item Z mno�iny $E_{i-1}$ spo��t�me mno�inu $E_i$ n�sledovn�:
$$
E_i =
\begin{cases}
E_{i-1} \cup \{e_i\} & \text{neobsahuje-li graf $(V,E_{i-1} \cup \{e_i\})$ kru�nici}\\
E_{i-1} & \text{jinak}
\end{cases}
$$
\end{penumerate}
Algoritmus se zastav�, pokud $E_i$ m� $n-1$ hran. Posledn� mno�ina $E_i$ jsou hrany minim�ln� kostry grafu $G.$
\end{e}

\begin{e}{V�ta}{0}{Slo�itost Kruskalova algoritmu}
P�i implementaci pot�ebujeme v podstat� vy�e�it �lohu udr�ov�n� ekvivalence (UNION-FIND): m�me mno�inu vrchol�, na po��tku je rozd�lena do jednoprvkov�ch t��d ekvivalence. Navrhn�te datov� struktury a algoritmus, kter� efektivn� vykon�v� dv� operace:
\begin{penumerate}
\item Sjednocen� t��d (UNION): u�init dva vrcholy ekvivalentn�mi t.j. nahradit t��dy je obsahuj�c� jejich sjednocen�m.
\item Testov�n� ekvivalence (FIND): Pro dan� dva vrcholy rozhodnout, zda jsou moment�ln� ekvivalentn�.
\end{penumerate}
V pr�b�hu algoritmu hled�n� minim�ln� kostry vykon�me $n-1$ operac� UNION --- jednu p�i ka�d�m p�id�n� hrany, a $m$ operac� FIND --- p�i ka�d�m testov�n�, zda p�id�van� hrana nevytvo�� kru�nici. \textit{(pozn.: Vrcholy t�to hrany mus� le�et v r�zn�ch komponent�ch.)}

�e�en�: vrcholy maj� p�i�azenu zna�ku ur�uj�c� t��du ekvivalence, do kter� pat��; a pro ka�dou t��du ekvivalence existuje seznam jej�ch vrchol�. Testov�n� ekvivalence je pak porovn�n� dvou zna�ek o slo�itosti $O(1)$ a p�i sjednocen� t��d mus�m p�ezna�it a p�em�stit vrcholy jedn� ze t��d, co� zabere $O(n)$ �asu. Pokud p�ezna�ujeme v�dy men�� t��du, vyjde odhad pot�ebn�ho �asu $O(n \log n + m).$
\end{e}

\begin{e}{Algoritmus}{0}{Jarn�k�v (Prim�v) algoritmus na hled�n� minim�ln� kostry grafu}
Je d�n souvisl� graf $G=(V,E).$ Budeme postupn� vytv��et mno�iny $V_0, V_1, \dots \subseteq V$ vrchol� a $E_0, E_1,\dots \subseteq$ hran, p�i�em� $E_0=0$ a $V_0=\{v\},$ kde $v$ je libovoln� zvolen� vrchol. V $i$-t�m kroku algoritmu vybereme z mno�iny hran $\{\{x,y\} \in E(G); x \in V_{i-1} \land y \in V \setminus V_{i-1}\}$ tu s minim�ln�m ohodnocen�m (bude to $e_i=\{x_i,y_i\}$) a polo��me $V_i=V_{i-1} \cup \{y_i\}$ a $E_i=E_{i-1} \cup \{e_i\}.$ Pokud ��dn� takov� hrana neexistuje, algoritmus kon��, to nastane v $(n-1)$-n�m kroku a $E_{n-1}$ je mno�ina hran minim�ln� kostry grafu. \textit{(Kostru tedy za��n�m budovat od jedin�ho vrcholu a v ka�d�m kroku k n� p�id�m nejkrat�� z hran mezi vrcholy kostry a zbytkem vrchol�.)} \footnote{Jarn�k�v algoritmus lze implementovat v �ase $O((n+m) \log n).$}
\end{e}

\begin{e}{Algoritmus}{0}{Topologick� t��d�n�}
\textbf{Probl�m:} V orientovan�m grafu $G=(V,E)$ sestrojte prost� zobrazen� $f: V \to \{1\dots|V(G)|\}$ tak, aby $\forall (v_1,v_2) \in E; f(v_1)<f(v_2).$ \textit{(Tedy m�me o��slovat vrcholy prvn�mi p�irozen�mi ��sly tak, aby hrany vedly jen z vrcholu s ni���m ��slem do vrcholu s vy���m ��slem.)}\hfil\break \textbf{Algoritmus:} Nejprve nastav�me ��ta� na 1. Nalezneme vrchol, do kter�ho nevede ��dn� hrana; tento vrchol o��slujeme ��ta�em a odtrhneme ho od grafu (sestroj�me indukovan� podrgraf) a zv���me ��ta�. Tento krok opakujeme, dokud mno�ina vrchol� grafu nen� pr�zdn�. Pokud nem��eme v n�jak�m kroku nal�zt bod, do kter�ho nevede ��dn� hrana, nelze graf topologicky set��dit.\hfil\break \textbf{Vyu�it�:} Odpov�daj�-li vrcholy grafu jednotliv�m krok�m n�jak�ho postupu a hrany �asov�m z�vislostem, kter� je t�eba zachovat, odpov�d� topologick� set��d�n� tohoto grafu (jednomu z) po�ad�, ve kter�m je nutn� kroky vykon�vat.
\end{e}


\end{document}
